\documentclass[12pt,a4paper,final]{article}
\usepackage[T1]{fontenc}
\usepackage[utf8]{inputenc}
\usepackage[english]{babel}
\usepackage{lmodern}
\usepackage{amsmath,amssymb,amsthm}
\usepackage[margin=2.5cm]{geometry}
\usepackage{hyperref}

\theoremstyle{plain}
\newtheorem{theorem}{Theorem}[section]

\author{Toufik Bellaj}
\title{HJM model density of the shifted lognormal case with stochastic volatility}
\date{}

\begin{document}
\maketitle

\section{Model formulation}
Starting from the multi-factor HJM term structure model for the dynamics of the
forward rates, driven by $N$ correlated Brownian motions, we have
\begin{equation}
df(t,T)
=\Biggl(\sum_{i,j}\sigma_i(t,T)\rho_{ij}(t)\int_{t}^{T}\sigma_j(t,u)\,du\Biggr)dt
+\sum_i \sigma_i(t,T)\,dW_i(t).
\end{equation}
Under the separability condition
\begin{equation}
\sigma_i(t,u)=\sigma_i(t,X)\,\theta_i(u)
\end{equation}
we obtain
\begin{eqnarray*}
df(t,T)
&=&\sum_{i,j}\theta_i(T)\sigma_i(t,X)\rho_{ij}(t)\sigma_j(t,X)
\Biggl(\int_{t}^{T}\theta_j(u)\,du\Biggr)dt
+\sum_i\theta_i(T)\sigma_i(t,X)\,dW_i(t)\\
&=&\sum_{i,j}\theta_i(T)\,d\Psi_{i,j}(t,X)\,\Theta_j(T)
+\sum_i\theta_i(T)\,dX_i(t),
\end{eqnarray*}
where
\begin{eqnarray}
\Theta_j(T)&=&\int_{t_0}^{T}\theta_j(u)\,du,\\
d\Psi_{i,j}(t,X)&=&\sigma_i(t,X)\rho_{ij}(t)\sigma_j(t,X)\,dt,\\
dX_i&=&\sigma_i(t,X)
\Biggl(dW_i(t)-\sum_j\rho_{ij}(t)\sigma_j(t,X)\Theta_j(t)\,dt\Biggr),
\end{eqnarray}
with $X_i(t_0)=0$ and $\Psi_{i,j}(t_0)=0$.

Integrating then yields
\begin{eqnarray*}
f(t,T)
&=& f(t_0,T)
+\sum_{i}\theta_i(T)
\Biggl(X_i(t)+\sum_{j}\Psi_{i,j}(t,X)\Theta_j(T)\Biggr),\\
r_t
&=& f(t_0,t)
+\sum_{i}\theta_i(t)
\Biggl(X_i(t)+\sum_{j}\Psi_{i,j}(t,X)\Theta_j(t)\Biggr)\\
&=& f(t_0,t)+\overline{r}_t,\\
P(t,T)
&=&\exp\Biggl(-\int_{t}^{T}f(t_0,u)\,du\Biggr),\\
B(t,T)
&=&P(t,T)Z_t\exp\Biggl(
-\sum_{i}\Theta_i(T)
\Biggl(X_i(t)+\frac12\sum_{j}\Psi_{i,j}(t,X)\Theta_j(T)\Biggr)
\Biggr),
\end{eqnarray*}
where
\begin{eqnarray*}
Z_t
&=&\exp\Biggl(
\sum_{i}\Theta_i(t)
\Biggl(X_i(t)+\frac12\sum_{j}\Psi_{i,j}(t,X)\Theta_j(t)\Biggr)
\Biggr).
\end{eqnarray*}
Introduce the Radon--Nikodym process $M_t$ by
\begin{eqnarray}
\frac{dM_t}{M_t}&=&\sum_i\Theta_i(t)\sigma_i(t)\,dW_i(t).
\end{eqnarray}

\begin{theorem}\label{thm:rn}
The Radon--Nikodym derivative $M_t$ satisfies
\begin{eqnarray*}
\frac{M_t}{M_s}
&=&\exp\Biggl(-\int_{s}^{t}\overline{r}_u\,du\Biggr)\frac{Z_t}{Z_s}.
\end{eqnarray*}
Under this measure the dynamics become
\begin{eqnarray*}
\Theta_j(T)&=&\int_{t_0}^{T}\theta_j(u)\,du,\\
d\Psi_{i,j}(t,X)&=&\sigma_i(t,X)\rho_{ij}(t)\sigma_j(t,X)\,dt,\\
dX_i&=&\sigma_i(t,X)\,dW^{M}_i(t),
\end{eqnarray*}
and
\begin{eqnarray*}
E_{t}^{Q}\Biggl[\exp\Biggl(-\int_{t}^{T} r_u\,du\Biggr)\Phi(X_T,\Psi_T)\Biggr]
&=&P(t,T)Z_t\,E_{t}^{M}\Biggl[\frac{\Phi(X_T,\Psi_T)}{Z_T}\Biggr].
\end{eqnarray*}
\end{theorem}

\begin{proof}
It\^{o}'s formula gives
\begin{eqnarray*}
\Theta_i(t)X_i(t)-\Theta_i(s)X_i(s)
&=&\int_{s}^{t}\theta_i(u)X_i(u)\,du
+\int_{s}^{t}\Theta_i(u)\,dX_i(u).
\end{eqnarray*}
Hence
\begin{eqnarray*}
\int_{s}^{t}\overline{r}_u\,du
&=&\sum_{i}\Biggl(
\Theta_i(t)X_i(t)-\Theta_i(s)X_i(s)
-\int_{s}^{t}\Theta_i(u)\sigma_i(u)\,dW_i(u)\Biggr)\\
&&+\sum_{i,j}\int_{s}^{t}\theta_i(u)\Psi_{i,j}(u)\Theta_j(u)\,du
+\sum_{i,j}\int_{s}^{t}\Theta_i(u)\sigma_i(u)\rho_{ij}(u)\sigma_j(u)\Theta_j(u)\,du.
\end{eqnarray*}
Integration by parts yields
\begin{eqnarray*}
\sum_{i,j}\int_{s}^{t}\theta_i(u)\Psi_{i,j}(u)\Theta_j(u)\,du
&=&\frac12\sum_{i,j}\Biggl(
\Theta_i(t)\Psi_{i,j}(t,X)\Theta_j(t)
-\Theta_i(s)\Psi_{i,j}(s,X)\Theta_j(s)\Biggr)\\
&&-\frac12\sum_{i,j}\int_{s}^{t}
\Theta_i(u)\sigma_i(u)\rho_{ij}(u)\sigma_j(u)\Theta_j(u)\,du.
\end{eqnarray*}
Therefore
\begin{eqnarray*}
\int_{s}^{t}\overline{r}_u\,du
&=&\sum_{i}\Theta_i(t)
\Biggl(X_i(t)+\frac12\sum_{j}\Psi_{i,j}(t,X)\Theta_j(t)\Biggr)\\
&&-\sum_{i}\Theta_i(s)
\Biggl(X_i(s)+\frac12\sum_{j}\Psi_{i,j}(s,X)\Theta_j(s)\Biggr)\\
&&-\sum_{i}\int_{s}^{t}\Theta_i(u)\sigma_i(u)\,dW_i(u)
+\frac12\sum_{i,j}\int_{s}^{t}
\Theta_i(u)\sigma_i(u)\rho_{ij}(u)\sigma_j(u)\Theta_j(u)\,du.
\end{eqnarray*}
The first two lines are $\log Z_t-\log Z_s$. The remaining stochastic
exponential is exactly $M_t/M_s$, which proves
\[
\frac{M_t}{M_s}
=\exp\Biggl(-\int_{s}^{t}\overline{r}_u\,du\Biggr)\frac{Z_t}{Z_s}.
\]
Under $M$, the finite-variation drift in $dX_i$ is removed by Girsanov,
and the pricing identity follows from $dQ|_{{\mathcal F}_T}/dM|_{{\mathcal F}_T}=M_T^{-1}$
together with $P(t,T)=\exp(-\int_t^T f(t_0,u)\,du)$.
\end{proof}

\section{Shifted lognormal case with lognormal stochastic volatility}
In this section we consider the one-factor case with volatility
\begin{equation}
\sigma(t,X)=\sigma(t)\,\eta_t\,\bigl(\beta (X + \gamma \Psi(t,X))+1\bigr),
\end{equation}
where $\eta_t$ follows an uncorrelated Black--Karasinski dynamics~\cite{BK}
\begin{eqnarray}
d\log\eta_t
&=&\kappa\bigl(m(t)-\log\eta_t\bigr)\,dt+\nu\,dZ_t,
\qquad
d\langle Z,W\rangle=0.
\end{eqnarray}
The zero correlation is used so that the time-changed Brownian motion is
independent of the clock and the Matsumoto--Yor formulae below apply. It is
not a material restriction on the smile. Instantaneous volatility is
$V_t=\eta_t S_t$, with $S_t=\beta(X+\gamma\Psi)+1$, so $dX$ and $dV$ remain
correlated through the local factor $S$ even though $\eta$ is independent of
$W$:
\begin{eqnarray*}
dX_t&=&\sigma(t)\,V_t\,dW^M_t,\\
dV_t&=&\beta\sigma(t)\,\eta_t V_t\,dW^M_t
+\nu\eta_t S_t\,dZ_t+\cdots,\\
\rho^{\mathrm{eff}}_t
&=&
\frac{\langle dX,dV\rangle}{\sqrt{\langle dX\rangle\langle dV\rangle}}
=
\frac{\beta\,\sigma(t)\,\eta_t}{\sqrt{\beta^{2}\sigma^{2}(t)\,\eta_t^{2}+\nu^{2}}}.
\end{eqnarray*}
Thus $\beta$ is the parameter that captures rate--vol correlation (the SABR
$\beta$--versus--$\rho$ trade-off). Once $\beta$ is free, a correlated driver
for $\eta$ is weakly identified and would only add unspanned skew at the cost
of destroying the time change. The remaining parameter $\gamma$ is not a
correlation: it mixes the quadratic-variation state $\Psi$ into $S$ and
produces the Girsanov kernel $(\gamma/\beta)S$ below.

Set
\begin{eqnarray*}
S_t&=& \beta (X +\gamma \Psi(t,X))+1,\\
A_t&=& \beta^2 \Psi(t,X).
\end{eqnarray*}
The initial conditions $X(t_0)=\Psi(t_0)=0$ give $S_{t_0}=1$ and $A_{t_0}=0$.
Using Theorem~\ref{thm:rn} under the $M$-measure we have
\begin{eqnarray}
\frac{dS_t}{S_t}
&=&\beta\sigma(t)\,\eta_t\,dW^M_t
+\beta\gamma\,\sigma^2(t)\,\eta_t^{2}\,S_t\,dt,\\
dA_t&=&\beta^2 \sigma^2(t)\,\eta_t^{2}\,S^2_t\,dt.
\end{eqnarray}

Define the stochastic time change
\begin{eqnarray*}
d\tau_t &=& \beta^2 \sigma^2(t)\,\eta_t^{2}\,dt,
\end{eqnarray*}
so that
\begin{equation}
\tau_t = \beta^2 \int_{t_0}^{t} \sigma^2(u)\,\eta_u^{2}\,du.
\end{equation}
Then, under $M$ and in the time-changed clock,
\begin{eqnarray}
\frac{dS_{\tau}}{S_{\tau}}&=&dW^{M}_{\tau}+\frac{\gamma}{\beta}\,S_{\tau}\,d\tau,\\
dA_{\tau}&=&S^{2}_{\tau}\,d\tau.
\end{eqnarray}

Girsanov with kernel $(\gamma/\beta)S$ removes the finite-variation drift.
The resulting density is
\begin{eqnarray*}
L_{\tau}
&:=&
\frac{dQ^{D}}{dQ^{M}}
=\exp\Biggl(-\frac{\gamma}{\beta}(S_{\tau}-1)+\frac{\gamma^{2}}{2\beta^{2}}A_{\tau}\Biggr),
\end{eqnarray*}
hence
\begin{eqnarray*}
\frac{dQ^{M}}{dQ^{D}}
&=&L_{\tau}^{-1}
=\exp\Biggl(\frac{\gamma}{\beta}(S_{\tau}-1)-\frac{\gamma^{2}}{2\beta^{2}}A_{\tau}\Biggr).
\end{eqnarray*}
Under $Q^{D}$ the dynamics become
\begin{eqnarray*}
S_{\tau}&=&\exp\Biggl(W^{D}_{\tau}-\frac{\tau}{2}\Biggr),\\
A_{\tau}&=&\int_0^{\tau} S^{2}_{u}\,du,\\
U_{\tau}&=&\frac{S_{\tau}}{A_{\tau}}.
\end{eqnarray*}
In particular $S_{\tau}=\exp(B_{\tau}^{(-1/2)})$ and
$A_{\tau}=A_{\tau}^{(-1/2)}$ in the notation of Matsumoto and Yor~\cite{YOR}.
Pricing identities under $M$ are rewritten under $D$ by
\begin{eqnarray*}
E^{M}\Biggl[\frac{g}{Z}\Biggr]
&=&
E^{D}\Biggl[L^{-1}\frac{g}{Z}\Biggr].
\end{eqnarray*}

The true one-factor HJM bond is, with $X=(S-1)/\beta-\gamma A/\beta^{2}$ and
$\Psi=A/\beta^{2}$,
\begin{eqnarray*}
\frac{B(t,T)}{Z_t}
&=&P(t,T)\exp\Biggl(
-\frac{\Theta(T)}{\beta}(S_{\tau}-1)
+\frac{\gamma\Theta(T)}{\beta^{2}}A_{\tau}
-\frac{\Theta(T)^{2}}{2\beta^{2}}A_{\tau}
\Biggr).
\end{eqnarray*}
Multiplying by the Radon--Nikodym factor $L^{-1}$ produces the $D$-measure
kernel
\begin{eqnarray*}
L_{\tau}^{-1}\frac{B(t,T)}{Z_t}
&=&P(t,T)\exp\Biggl(
-\Biggl(\frac{\Theta(T)-\gamma}{\beta}\Biggr)(S_{\tau}-1)
-\frac{A_{\tau}}{2}
\Biggl(\frac{\Theta(T)-\gamma}{\beta}\Biggr)^{2}
\Biggr).
\end{eqnarray*}
(The previous display is \emph{not} the bond $B/Z$; it is the integrand for
an expectation under $Q^{D}$.)

Write $\alpha(T)=(\Theta(T)-\gamma)/\beta$. The joint law of
$(S_{\tau},U_{\tau})$ under $Q^{D}$, conditionally on the clock $\tau$ (i.e.\
on the path of $\eta$), follows from~\cite[Thm.~4.1]{YOR} after the change of
variables $U=S/A$. The conditional law of $S$ given $U=u$ is inverse Gaussian
with mean $1$ and shape $u$,
\begin{eqnarray*}
p(S_{\tau}=x\mid U_{\tau}=u,\tau)
&=&\sqrt{\frac{u}{2\pi x^{3}}}
\exp\Biggl(u-\frac{u}{2}\Biggl(x+\frac{1}{x}\Biggr)\Biggr),
\end{eqnarray*}
and the marginal of $U$ is
\begin{eqnarray*}
p(U_{\tau}=u\mid\tau)
&=&\sqrt{2\pi}\,\exp\Biggl(-u-\frac{\tau}{8}\Biggr)
\frac{H(u,\tau)}{u^{3/2}},
\end{eqnarray*}
where $H$ is the Hartman--Watson function $\theta$ of~\cite{YOR}. The
conditional density in $S$ does not depend on $\tau$.

The product of this density with the $D$-kernel, divided by $P(t,T)$, is the
density of an inverse Gaussian law with mean $1/\mu$ and shape $U_{\tau}$,
where $\mu=1+\alpha(T)/U_{\tau}=1+(\Theta(T)-\gamma)/(\beta U_{\tau})$:
\begin{eqnarray*}
p(x\mid U_{\tau},\tau)\,
L^{-1}\frac{B(t,T)}{P(t,T)Z_t}\Biggm|_{S=x}
&=&\sqrt{\frac{U_{\tau}}{2\pi x^{3}}}
\exp\Biggl(
\alpha(T)+U_{\tau}
-\frac12\Biggl(
\frac{x}{U_{\tau}}
\bigl(U_{\tau}+\alpha(T)\bigr)^{2}
+\frac{U_{\tau}}{x}
\Biggr)
\Biggr)\\
&=&\sqrt{\frac{U_{\tau}}{2\pi x^{3}}}
\exp\Biggl(
-\frac{U_{\tau}}{2x}
\Biggl(x\mu-1\Biggr)^{2}
\Biggr)\\
&=&\sqrt{\frac{U_{\tau}}{2\pi x^{3}}}
\exp\Biggl(
-\frac12\Biggl(\sqrt{U_{\tau}x}\,\mu-\sqrt{\frac{U_{\tau}}{x}}\Biggr)^{2}
\Biggr)\\
&=&\sqrt{\frac{U_{\tau}}{2\pi x^{3}}}
\exp\Biggl(
2\mu U_{\tau}
-\frac12\Biggl(\sqrt{U_{\tau}x}\,\mu+\sqrt{\frac{U_{\tau}}{x}}\Biggr)^{2}
\Biggr).
\end{eqnarray*}
(The four right-hand sides are identical; they are a density times a kernel,
not the bond itself.)

Define
\begin{eqnarray*}
\xi^{+}(k)&=&\sqrt{U_{\tau}k}\,\mu-\sqrt{\frac{U_{\tau}}{k}},\\
\xi^{-}(k)&=&-\sqrt{U_{\tau}k}\,\mu-\sqrt{\frac{U_{\tau}}{k}}.
\end{eqnarray*}
At frozen $U_{\tau}$ one has $d\xi^{+}+d\xi^{-}=\sqrt{U_{\tau}/k^{3}}\,dk$,
but the two branches carry different Jacobians; the Gaussian split uses each
branch separately. Integrating the product against $\mathbf{1}_{\{S\le k\}}$
yields the inverse-Gaussian cdf (mean $1/\mu$, shape $U_{\tau}$)
\begin{eqnarray*}
E^{D}\Biggl[
L^{-1}\frac{B(t,T)}{Z_t}
\mathbf{1}_{\{S_{\tau}\le k\}}
\Biggm|U_{\tau},\tau\Biggr]
&=&P(t,T)
\Biggl(
N\bigl(\xi^{+}(k)\bigr)
+\exp\Biggl(2\Biggl(U_{\tau}+\frac{\Theta(T)-\gamma}{\beta}\Biggr)\Biggr)
N\bigl(\xi^{-}(k)\bigr)
\Biggr).
\end{eqnarray*}

A payer swaption with expiry $T_{e}$, last payment $T_{n}$ and strike $K$ has
time-$T_{e}$ value
$(1-B(T_{e},T_{n})-K\sum_{i=e+1}^{n}\delta_{i}B(T_{e},T_{i}))^{+}$,
using $B(T_{e},T_{e})=1$. Theorem~\ref{thm:rn} gives
\begin{eqnarray*}
\mathrm{swap}(t)
&=&P(t,T_{e})Z_{t}\,
E^{M}\Biggl[
\frac{
\bigl(1-B(T_{e},T_{n})-K\sum_{i=e+1}^{n}\delta_{i}B(T_{e},T_{i})\bigr)^{+}
}{Z_{T_{e}}}
\Biggr]\\
&=&P(t,T_{e})Z_{t}\,
E^{D}\Biggl[
L^{-1}
\bigl(G_{e}-G_{n}-K\sum_{i=e+1}^{n}\delta_{i}G_{i}\bigr)^{+}
\Biggr],
\end{eqnarray*}
where $G_{i}:=B(T_{e},T_{i})/Z_{T_{e}}$. Given $U$, every $G_{i}$ is strictly
monotone in $S$, so the forward swap value has a unique root $S^{*}$ and the
positive part is the same linear combination restricted to $\{S>S^{*}\}$ (or
$\{S<S^{*}\}$, according to the sign of $\beta$). Each term
$E^{D}[L^{-1}G_{i}\mathbf{1}_{\{S\gtrless S^{*}\}}\mid U,\tau]$ is the
inverse-Gaussian formula above with tenor-specific
$\alpha(T_{i})=(\Theta(T_{i})-\gamma)/\beta$. The remaining average is over
the Hartman--Watson law of $U$ given $\tau$ and over the law of the clock
$\tau=\beta^{2}\int_{t_0}^{T_{e}}\sigma^{2}(u)\eta_{u}^{2}\,du$. Under
Black--Karasinski, $\log\eta$ is Ornstein--Uhlenbeck, so this mixing law is
the exponential functional $\int e^{2\log\eta_u}\,du$ of an OU
process~\cite{YOR}.

\begin{thebibliography}{9}
\bibitem{YOR}
H.~Matsumoto and M.~Yor,
Exponential functionals of Brownian motion, I: Probability laws at fixed time,
\textit{Probability Surveys} \textbf{2} (2005), 312--347.

\bibitem{BK}
F.~Black and P.~Karasinski,
Bond and option pricing when short rates are lognormal,
\textit{Financial Analysts Journal} \textbf{47} (1991), 52--59.

\bibitem{lognormal_vol}
L.~B.~G.~Andersen and V.~V.~Piterbarg,
\textit{Interest Rate Modeling. Volume~I: Foundations and Vanilla Models},
Atlantic Financial Press, 2010.
\end{thebibliography}

\end{document}
